\section{绪论}

\subsection{研究背景与意义}

近年来，大语言模型在自然语言理解与生成领域迅速发展。以 Transformer 为基础的生成式模型不断扩大参数规模，在问答、摘要、翻译、代码补全和智能助手等任务中表现出较强的任务处理能力\cite{vaswani2017attention,touvron2023llama2}。与此同时，模型推理成本也持续上升。大模型生成通常采用自回归方式逐词元输出，推理过程具有串行性，因此即使单步计算经过工程优化，整体延迟仍然偏高\cite{leviathan2023speculative,kwon2023pagedattention}。

对于移动端场景而言，这一矛盾尤为突出。移动终端通常承担人机交互入口，但其处理器性能、内存空间和功耗预算都低于服务器或桌面工作站。量化、小模型部署以及高效推理框架能够缓解本地部署困难\cite{ggerganov2025llamacpp,frantar2022gptq,alizadeh2023llmflash}，但移动端若要稳定完成高质量大模型推理，仍然面临性能和资源限制。

云端或桌面端计算能力较强，适合运行参数规模更大、生成质量更高的模型。然而纯远程推理同样存在问题：网络请求和响应会引入额外时延，局域网或互联网环境中的带宽波动会放大这种不确定性，短文本或多轮对话场景下的交互等待尤为明显。若全部生成都在远端完成，则移动端本地算力无法参与实际推理过程，端侧计算资源也难以转化为生成阶段的有效吞吐。

端云协同推理的意义在于把移动端的即时交互能力和远端的高质量模型能力结合起来，使终端从输入显示入口扩展为生成过程中的轻量计算节点。对于大语言模型而言，投机解码适合作为这种协同的切入点\cite{leviathan2023speculative,chen2023speculative,miao2023specinfer,li2024eagle}：草稿侧以较低计算开销提出候选词元，目标侧对候选进行验证。候选被接受时，系统一次推进多个词元；候选被拒绝时，目标模型给出修正结果继续生成。该机制把逐词元等待的目标模型调用转化为“草稿提案 + 批量验证”的协同过程，有利于降低用户感知等待时间。

现有投机解码研究多以单机或服务器环境为对象，已经形成草稿模型、词元树、多头候选和检索辅助等多种方案\cite{leviathan2023speculative,miao2023specinfer,cai2024medusa,he2023rest}。在端云协同场景中，草稿侧和目标侧位于不同设备，系统还需要处理跨设备状态同步、会话管理、词元空间一致性以及验证开销控制。基于上述背景，本文构建面向 Android 移动端与桌面端协同场景的投机解码实验系统，并通过实验分析其运行效果与瓶颈。

\subsection{国内外研究现状}

围绕大语言模型推理效率问题，已有研究从模型压缩、量化部署、缓存管理、并行解码和服务系统优化等多个方向展开\cite{ggerganov2025llamacpp,kwon2023pagedattention,frantar2022gptq,dao2022flashattention,song2023powerinfer}。其中，面向推理加速的工作大致可分为两类：一类关注降低单次前向计算成本，例如权重量化、KV Cache 优化、注意力计算优化和内存调度；另一类关注减少高质量目标模型实际参与解码的次数，投机解码属于这一类。

Leviathan 等提出的 Speculative Decoding 方法给出了草稿模型和目标模型协同工作的基本框架，在一定条件下通过草稿提案和目标验证减少整体解码时间\cite{leviathan2023speculative}。Chen 等讨论了 speculative sampling 的实现思路，强调草稿生成与目标校正之间的概率关系\cite{chen2023speculative}。Miao 等提出的 SpecInfer 从系统角度讨论词元树验证，使目标模型在一次验证中处理多个候选分支\cite{miao2023specinfer}。Medusa 通过在模型上增加多个解码头生成候选，减少外部草稿模型依赖\cite{cai2024medusa}；Lookahead Decoding 试图打破传统自回归推理中的部分顺序依赖\cite{fu2024lookahead}；REST 则利用检索结果构造候选，从而减少草稿模型计算开销\cite{he2023rest}。Li 等提出的 EAGLE 强调统一词元空间、条件概率一致性以及输出分布保持问题\cite{li2024eagle}。这些研究表明，投机解码的效果同时受候选生成速度、候选质量、验证方式和状态一致性影响。

端云协同推理方面，现有研究多从任务卸载、边缘智能、移动 AI 部署和分层推理等角度展开\cite{satyanarayanan2017edge,shi2016edge,alizadeh2023llmflash}。在实际系统中，移动端通常作为输入和显示入口，远端承担主要推理任务。这种模式易于实现，但终端参与度有限，多数系统仍采用“移动端提交请求、远端完整生成、移动端展示结果”的传统方式。本文关注的方案使移动端参与草稿生成，并与远端目标模型共同构成投机解码过程。

现有研究仍存在两个不足。其一，许多工作更关注单机环境中的理论与算法正确性，而对跨设备场景下的协议设计、会话状态和工程边界关注不足。其二，许多研究对词元对齐和运行时连续性问题的讨论仍停留在算法接口层，而在真实移动端与桌面端协同条件下，这些问题会直接影响系统运行稳定性。基于此，本文将研究重点放在端云协同投机解码系统的设计、实现和实验分析上。

\subsection{本文研究内容}

本文围绕 Android 移动端与桌面端协同场景下的投机解码展开研究，主要内容包括以下几个方面。

首先，设计了“移动端草稿生成、桌面端目标验证”的投机解码总体方案。Android 端负责草稿候选生成，桌面端推理服务负责目标验证，双方围绕统一的投机会话开展多轮协同。

其次，构建了端云协同投机解码协议与会话管理机制。系统实现了 \texttt{start}、\texttt{propose}、\texttt{fallback} 和 \texttt{close} 等基础消息，并设计了草稿会话与目标会话两类状态对象，用于维护已接受前缀、候选提案和目标验证过程。

再次，围绕词元对齐与验证真实性问题，系统采用统一真实词元空间下的目标验证路径，支持线性提案与树形候选验证，并结合运行时连续性优化降低协同过程中的中间开销。

最后，本文对所实现系统进行了实验验证。实验中将云侧本地单进程投机加速和云侧本地双进程投机加速作为主要对照，重点分析跨设备投机解码系统的完整流程能力、接受率、吞吐率以及验证侧开销变化。

\subsection{论文结构安排}

本文共分为六章。第一章介绍研究背景、国内外研究现状和本文的主要研究内容。第二章介绍大语言模型推理、端云协同和投机解码的相关技术基础。第三章给出端云协同投机解码方案设计，包括系统结构、协议与关键问题分析。第四章详细介绍投机解码系统的实现。第五章给出实验设计、实验结果及分析。第六章对全文进行总结，并讨论系统的不足与未来改进方向。
